\documentclass[11pt]{article}
\usepackage[top=1.00in, bottom=1.0in, left=1.1in, right=1.1in]{geometry}
\usepackage{graphicx}
\usepackage{natbib}
\usepackage{amsmath}
\usepackage{lineno}
\usepackage{gensymb}
\usepackage{xr-hyper}
\usepackage{hyperref}
\setlength\parindent{0pt}


\begin{document}
\setlength{\parindent}{0cm}
\setlength{\parskip}{7pt}


\renewcommand{\refname}{\CHead{}}

Editor and reviewer comments (we provide below the full context of the two reviewers' comments) are in \emph{italics}, while our responses are in regular text. \\ 

{\bf Response to editor's comments:} 

\emph{Both reviewers found this contribution quite interesting. Rev. 1 raises an important question about the overlooked role of moisture availability as an additional, interacting and important control over phenology. This and all other points raised by both reviewers should be thoroughly addressed in a revision. In addition, I have two more comments.}

We appreciate the opportunity to revise our manuscript. We have worked to clarify the points raised by the 2 reviewers in this new version.

\emph{First, and in contrast to Rev. 2, I found the hypothesis(es) hard to disentangle. The Introduction begins by discussing how plants use cues to "adjust" their phenology to variable conditions, but then segues quickly into plants' use of the fixed cue of summer solstice, without explaining this seeming disconnect. It then begins referring to "this hypothesis" and "this proposed photoperiod switch" (and so on) without ever clearly and simply stating the hypothesis. In the results, the "hypothesis that solstice is not a reliable cue" is then proposed, which violates the standard definition of a hypothesis in that it states a lack of effect rather than an effect. Hypotheses stated in more definite wording would greatly help the flow of the paper}

[...]\\

\emph{Second, and as mentioned by Rev. 2, I wondered about the implications of the particular geographic (including latitudinal) scope of this Europe-based study. For an idea that is being put forward as a general explanation of a global phenomenon, more attention needs to be paid (at least verbally) to its applicability to the rest of the world.}

[...]\\

{\bf Reviewer comments:} 

{\bf Reviewer 1}

\emph{This is an interesting assessment of the potential importance of summer solstice in cueing phenophase timing in terrestrial plants. The analysis is largely sound, but the sole focus on growing degree days (GDD) in defining growth potential beyond a given day of year is a little concerning.}

We thank the reviewer for their positive comments regarding our manuscript and appreciate that they see the value in these analyses.\\

\emph{In particular, there is no consideration given to potential tradeoffs involving moisture availability. Garonna et al (2018 Environ. Res. Lett. 13 024025) argued pretty convincingly for shifting tradeoffs among photoperiod, temperature, and vapor pressure deficit as primary constraints on seasonal growth timing, with the implication being that alleviations or strengthening of one or two of those factors will lead to strengthening or alleviation, respectively, of the third. So while the focus on GDD is understandable, is it reasonable to expect the importance of heat accumulation to remain constant if it's also accompanied by increasing moisture limitation? The approach/framework feels a little like begging the question. I'm not suggesting that makes the manuscript uninteresting or unimportant, but it does leave me uncertain of the fit for PNAS.}

[...]\\

\emph{Minor comments:\\
Lines 4-5, the authors state "the fitness landscape shaping selection on these cues remain largely unknown for many events".}

[...]\\

\emph{This is a little vague - if it's important and a guiding motivation for this paper then provide an example. Also, "remain" should be "remains".}

[...]\\

\emph{Lines 7-8, the authors state "Recently, summer solstice has been proposed as a universal trigger to modulate cues and initiate key physiological processes".}

[...]\\

\emph{Here I would suggest adding a little more information such as, for example, in what kinds of plants or systems?}

[...]\\

\emph{Lines 15-17, the authors state "This proposed photoperiod switch, if correct, could reshape predictions of forest responses to climate change. Using a fixed date like the solstice as a cue, however, could limit plasticity in how plants respond across their ranges, which span very different climates."}

[...]\\

\emph{First, what is meant by photoperiod switch? This needs a little clarification. Second, this statement might be tempered by insights from Garonna et al. 2018 as mentioned above.}

[...]\\

\emph{Line 28: "exposition" is not the right word here. Consider changing to "risk of exposure".}

[...]\\

\emph{Line 33: "has" should be "have"}

[...]\\

\emph{Lines 36-37, the authors state "This heat accumulation is a key factor in development and growth processes of both crops and wild plants."}

[...]\\

\emph{Not arguing with that claim, but its importance and predictive value varies with other constraints such as moisture availability, as should be noted.}

[...]\\

\emph{Figure 1 legend: "stand" in the third line should be "stands". Also, in the final sentence it should be stated that moisture availability is assumed to remain constant if that is indeed the case.}

[...]\\

\emph{"Our results suggest solstice could act as a reliable marker but also highlight the challenges in disentangling the influence of the solstice from that of a thermal optimum cue."}

[...]\\

{\bf Reviewer 2}

\emph{Overall, this is a very interesting study in which the authors explore temperature accumulation to the summer solstice is a cue for plant growth transitions and potential mechanisms that could explain this trigger point. They test if temperature accumulation up to this point relative to remaining \# of GDDs is optimal. The results suggest more research is needed on thermal cues and perhaps interaction with photoperiod. The manuscript does an excellent job at synthesizing a complex topic and presenting it in an easy to understand manner. It is well written and presented. All figures and tables are relevant, and the statistical analyses are appropriate.\\
 I recommend this manuscript for publication as it builds on recent hypotheses around the summer solstice as an environmental cue for plant development while providing alternative suggestions, further discussion points and potential future research.}

[...]\\

\emph{Abstract\\
In my opinion the abstract is too short and lacks key information such as the nature of the data being explored, a time frame, and a rough geographic region....mainly to quantify what is meant by 'local'. It is also not clear what is meant by 'juncture'.}

[...]\\

\emph{Introduction \\
First 2-3 paragraphs: It would be useful to mention the general location of the studies, I presume it is temperate latitudes?}

[...]\\

\emph{L17 - leaf unfolding of what?}

[...]\\

\emph{Last 2 paragraphs: I agree that AGDDs is an important factor for phenological/growth progression but I think chilling requirements need a mention. Also GDD should be plural - i.e. GDDs - in most cases.}

[...]\\

\emph{L43-48 consider replacing the first instance of 'GDD' with 'accumulated temperature'. I do not really like the word predict - but can't think of a better word or phrase.}

[...]\\

\emph{L46 and 47 consider replacing 'enough GDD' with 'sufficient temperature'}

[...]\\

\emph{L49 consider replacing the first instance of 'GDD' with 'temperature'}

[...]\\

\emph{L50 is the base temperature 0C or some other value?}

[...]\\

\emph{Results and discussion\\
Figure 1: What does the vertical dashed line indicate in b, c and d? Summer solstice? This figure is excellent and very convincing.}

[...]\\

\emph{L69 consider replacing 'only' with 'on' or some other rewording of this sentence to make more sense.}

[...]\\

\emph{L71-78: excellent consideration of the possibility of coincidence!}

[...]\\

\emph{Consider using 'summer solstice' throughout.}

[...]\\


\end{document}
